\scanpage{218}\mspage{\handwritten{-4-}}
cation $\handwritten{\underline P'^{\underline q_{\underline{\unsure{o|0}}}}/\text{\textcircled{$\mathfrak m$}}_{\underline A}\underline P'^{\underline q_{\underline o}}\longrightarrow\underline P^{\underline q_{\underline o}}/\text{\textcircled{$\mathfrak m$}}_{\underline A}\underline P^{\underline q_{\underline o}}}$ étant injective,il en résulte que le complexe
\[
\handwritten{\cdots\longrightarrow\underline P^{\underline q_{\underline o}-2}\longrightarrow\underline P^{\underline q_{\underline o}-1}\longrightarrow\underline P'^{\underline q_{\underline o}}\longrightarrow0\longrightarrow\cdots}\ .
\]
vérifie les conditions du lemme. On peut donc supposer dans l'énoncé que $\handwritten{\underline P^{\underline q}=0}$ pour $\underline q>\underline q_{\underline o}$.Utilisons [(\handwritten{1}), Ch.III,§.5] pour montrer que $\underline H$ est $\underline A$-plat.Le complexe $\underline P^\bullet$ est une résolution pour $\underline H$ par des $\underline A$-modules plats.Par consequence,les modules
$\handwritten{\operatorname{Tor}^{\underline A}_\bullet(\underline A/\text{\textcircled{$\mathfrak m$}}_{\underline A},\underline H)}$ sont les modules de cohomologie du complexe $\handwritten{\underline P^{\unsure{\circ|\bullet}}\otimes_{\underline A}\underline A/\text{\textcircled{$\mathfrak m$}}_{\underline A}=\underline P^\circ/\text{\textcircled{$\mathfrak m$}}_{\underline A}\underline P^\circ}$ et la conclusion résulte.

\uline{Corollaire}. \uline{Soit $\underline f:\underline X\handwritten{\longrightarrow}\underline Y$ un morphisme propre et plat de variétés algébriques. Alors l'ensemble}
\[
\underline{\handwritten{\underline G_{\underline f}(\underline Y)=\{\underline y\in\underline Y\mid\underline X_{\underline y}}\quad\text{soit de Gorenstein}\ \handwritten{\}}}
\]
\uline{est ouvert dans $\underline Y$.}

\begin{flushright}L'Institut Central de Mathématique\\
14,\deleted{\unsure{x|r}}rue de l'Académie , Bucarest\\
\handwritten{ROUMANIE}\end{flushright}

\noindent\begin{minipage}[t]{.33\linewidth}
\diagram{\(\begin{array}{cl|}\text{---}&\handwritten{\text{latine}}\\
\handwritten{\text{\tikz[baseline=-.4ex]\draw[densely dotted] (0,0) circle (.63em);}}&\handwritten{\text{gothique}}\\
\handwritten{\bigcirc}&\handwritten{\text{grec}}\\
\handwritten{\text{\tikz[baseline=-.4ex]\draw[black!25] (0,0) circle (.63em);}}&\handwritten{\text{ronde}}\\ \hline
\end{array}\)}
\end{minipage}\hfill\begin{minipage}[t]{.63\linewidth}
\noindent\rule{\linewidth}{.3pt}\par
\handwritten{\uline{Bibliographie}}

\handwritten{[1] Bourbaki, N. -- Algèbre commutative, Paris, Hermann 1961}

\handwritten{[2] Grothendieck, A -- E.G.A., IV , 1965}

\handwritten{[3] Samuel, P -- Séminaire d'Algèbre commutative, 1966/67.}

\handwritten{[4] J. P. Serre -- Faisceaux algébriques cohérents, Ann. Math. 1955}
\end{minipage}
% archive sheet number: 217
